\section{对象与运动：稳定子、吸引域与激活轨迹}

\subsection{全局演化算子}

设全局系统状态为 $\{\mathrm{State}_t(v)\}_{v\in V}$。定义演化算子
$$
\mathsf T:\{\mathrm{State}_t(v)\}_{v\in V}\mapsto \{\mathrm{State}_{t+1}(v)\}_{v\in V},
$$
其中每个节点更新满足有限邻域性，并包含：
\begin{enumerate}
  \item 局部读出与折叠 $\Fold_{m(v)}$；
  \item 残差更新与连接提升 $\mathcal P$；
  \item 误差驱动的分辨率调度（octree/AMR）。
\end{enumerate}

\subsection{对象定义：周期轨与吸引域}

取有限节点集合 $S\subset V$。若存在周期 $p\ge 1$ 使
$$
\mathsf T^{p}\bigl(\mathrm{State}(S)\bigr)=\mathrm{State}(S),
$$
并且对小扰动 $\delta\mathrm{State}(S)$ 存在吸引域使轨道收敛回该周期轨，则称 $(S,p)$ 定义对象（稳定子）。该定义将“对象”归结为离散动力系统的稳定结构，不要求材料恒定，仅要求模式稳定。

\subsection{运动的关系化表述}

定义对象的激活位置序列为 $\{v_t\}\subset V$，其中 $v_t$ 为在时刻 $t$ 上对象模式的主要支撑节点集合（可取为能量/误差/残差密度最大的节点或其连通分量）。则“运动”可被表述为激活序列在 $(V,\mathrm{Nbr})$ 上的迁移。该迁移受两类约束控制：

\begin{enumerate}
  \item 拓扑约束：只能在有限邻域内传播（由 $\mathrm{Nbr}$ 决定）；
  \item 时间纤维约束：迁移必须伴随残差纤维上的一致性提升（由 $\mathcal P$ 与 holonomy 限制）。
\end{enumerate}

\subsection{\texorpdfstring{与 $C_m$ 与 $B_m$ 的结构联系（预告）}{与 Cm 与 Bm 的结构联系（预告）}}

在最小核 $m=6$ 中，循环语言 $C_6$ 的 $18$ 个元素可作为“无边界缺陷”的局部模式候选，而边界集合 $B_6$ 的 $3$ 个元素可作为“端点缺陷模式”候选。对象稳定性与运动可在后续章节中进一步与“缺陷传播”联系：缺陷在含接口回路族上诱导非平凡 holonomy，从而提高局部时间成本并触发分辨率调度。这一机制为“对象—缺陷—时间成本—分辨率”的闭环结构提供统一描述，并为第 13 节双边界扩展与第 15 节实验提供结构接口。
